% Research fragment. Source statements checked against the incoming files
% Alternating_Tails_Spectral_Defects and Quartic_Spectral_Reconstruction.
% Suggested bibliography keys: tail-alternating, tail-quartic.
\section{Complete optimal reconstruction at smaller tail order ten}
\label{t10:section}

Two recent source theorems leave only one pair at smaller order ten.
The order-ten theorem of \emph{Alternating Tails and Spectral Defects},
Theorem~\texttt{rt:rank10}, proves the conjecture for every integer
$b\ge6$ outside the progression $b\equiv962\pmod{2048}$.
The effective theorem of \emph{Quartic Spectral Reconstruction},
Theorem~\texttt{ert:main} and its exact threshold calculation,
proves it for every $b\ge1920$. We certify the single remaining
value $b=962$ by a five-square determinant over a finite field.
This completes the whole smaller-order-ten family. The method does
not assert that the unrestricted unequal-order conjecture is proved.

\subsection{The formal coordinate problem}

For even $p,q$ put
\[
 A_p(n)=\zeta(p,n+1)=\zeta(p)-H_n^{(p)},\qquad x_n=n+\tfrac12,
\]
and initially define
\[
 D_{p,q}(s)=\sum_{n\ge0}A_p(n)A_q(n)x_n^{-s},
 \qquad \Re s>3-p-q.
\]
Let $D_r=D_{p,q}(-2r)$ denote its regular continued centered moments.
The centered Hurwitz expansion has powers $x_n^{1-p-2j}$; hence
the product has only even inverse powers. Its possible continuation
poles lie at odd integers, so all the displayed $-2r$ values are regular.
The inherited centered continuation and finite Bernoulli reduction
supply inverse transforms
\begin{align}
 C_j&=\sum_{r=0}^j\binom{2j+1}{2r}2^{-2(j-r)}D_r,
 \label{t10:inverse}\\
 D_r&=\sum_{j=0}^r\frac{\binom{2r}{2j}}{2j+1}
 B_{2r-2j}(1/2)C_j.\label{t10:forward}
\end{align}
Write $p=2a$, $q=2b$, $a<b$. To specify the algebraic claim completely,
let $\mathcal V_{a,b}$ be the $2b$-dimensional rational vector space with
\emph{formal basis}
\[
 C_0,\ldots,C_{b-1},Z_0,\ldots,Z_{b-1}.
\]
These are independent labels by definition. Under numerical evaluation,
$Z_k$ is assigned $\zeta(2k)$, including $Z_0=-1/2$, and each early
$C_j$ is assigned its continued value above. No injectivity of this
evaluation map is asserted. For every $j\ge b$, the source reduction
defines $C_j=\sum_{k=0}^{b-1}M_{j-b,k}Z_k$, using the formula below
also for row indices greater than $b-1$. Thus every later $D_r$ is a
specified vector of $\mathcal V_{a,b}$. The coordinate matrix of
$C_0,\ldots,C_{2b-1}$ is $\operatorname{diag}(I_b,M_{a,b})$.
It follows that the first $2b$ moment coordinates have full formal
rank precisely when the matrix
\begin{equation}\label{t10:M}
 (M_{a,b})_{ik}=B_{2(i+k-a+1)}
 \left\{\frac{\binom{2i+2}{2a-2k}}{2i+2}
 +\frac{\binom{2b+2i-2a+2}{2b-2k}}{2b+2i-2a+2}\right\},
 \qquad0\le i,k<b,
\end{equation}
is nonsingular. A negative Bernoulli index gives zero. A binomial
summand is omitted unless its lower index is at least two and no
greater than its upper index. Thus a lower index zero is omitted.
The complete first-$2b$ coordinate determinant is
$\det M_{a,b}/(4b-1)!!$. These are inherited identities, not new
independence assertions about numerical zeta values.

\begin{theorem}[Complete smaller-order-ten reconstruction]
\label{t10:complete}
For every integer $b\ge6$, $M_{5,b}$ is nonsingular. Consequently,
for every even $q>10$, the moments
\[
 D_{10,q}(0),D_{10,q}(-2),\ldots,D_{10,q}(-2q+2)
\]
have formal coordinate rank $q$ in the inherited coordinate space.
Every later centered moment is a finite rational linear combination
of these first $q$ moments.
\end{theorem}

The only new nonvanishing certificate needed by the theorem is
\begin{equation}\label{t10:main-certificate}
 \boxed{\det M_{5,962}\equiv15578\pmod{65537}.}
\end{equation}
The rational determinant has denominator prime to $65537$;
therefore this congruence proves exact nonvanishing.

\subsection{The retained five-square matrix}

We state the inherited contraction explicitly to make the finite
certificate independently reproducible. For $b\ge2a$ and
$0\le i,\ell<a$, define
\[
 n_i=b-a+i,\quad N_{i\ell}=2(n_i-\ell)+1,\quad
 R_\ell=2b-2a+2+2\ell,
\]
\[
 F_{i\ell}=\frac1{(2\ell+1)!}
 \prod_{r=0}^{2\ell-1}(2n_i-r)(2n_i+3+r).
\]
Define the integer Newton coefficients $d_{ih}$ by
\begin{equation}\label{t10:Newton}
 (2n_i+2-y)(4n_i+2-y)^{\underline{2i}}
       =\sum_{h=0}^{2i+1}d_{ih}y^{\underline h}.
\end{equation}
Here $z^{\underline h}=z(z-1)\cdots(z-h+1)$. Equivalently,
if the polynomial on the left is $W_i(y)$, then
\[
 d_{ih}=\frac1{h!}\sum_{j=0}^h(-1)^{h-j}\binom hj W_i(j).
\]
The retained matrix is
\begin{equation}\label{t10:A}
 \mathcal A_{i\ell}=2F_{i\ell}
 +\frac{R_\ell!}
 {(2n_i+1)(2n_i+2)(2\ell+2)!N_{i\ell}!}
 \sum_{k=0}^i d_{i,2k+1}N_{i\ell}^{\underline{2k+1}}
                 2B_{2b-2(a-i+\ell+k)}.
\end{equation}
The omitted lower-index-zero convention in \eqref{t10:M} is
already included in the $2F$ term and must not be altered.

For completeness, the source multiplier is
\begin{equation}\label{t10:multiplier}
 \begin{aligned}
 \det M_{a,b}={}&(-1)^{a(b-a)+a(a-1)/2}
 \frac{\prod_{t=0}^{b-1}(2b+2t-2a+1)!}
      {\prod_{k=0}^{b-1}(2b-2k)!}\\
 &\quad\cdot\prod_{n=0}^{b-a-1}\frac{((2n+1)!)^2}{(4n+2)!(4n+3)!}
 \prod_{i=0}^{a-1}\frac{(2n_i+1)(2n_i+2)!}{(4n_i+2)!}
 \det\mathcal A.
 \end{aligned}
\end{equation}
The exact derivation is the Pad\'e-row elimination in the source's
Lemma~\texttt{ert:normalization}. It can also be replayed directly:
the monic rows
\[
 p_{nt}=\frac{2(2n+2t+1)!(2n+1)!}
 {(4n+2)!(2n-2t+1)!(2t)!},\qquad0\le t\le n,
\]
are applied to the factorial-normalized Bernoulli matrix. Their
orthogonality eliminates the last $b-a$ columns from the final $a$
rows. Reversing the retained columns, and dividing row $i$ by
$c_{n_i}=(2n_i+1)(2n_i+2)!/(4n_i+2)!$, gives \eqref{t10:A}.
The diagonal eliminated entries are
$((2n+1)!)^2/((4n+2)!(4n+3)!)$, which gives
\eqref{t10:multiplier}, including its sign.

\subsection{The complete finite certificate}

Set $a=5$, $b=962$, and $P=65537$. The prime $P$ exceeds every
factorial argument in \eqref{t10:A} and \eqref{t10:multiplier}.
The Bernoulli recurrence
\begin{equation}\label{t10:Brec}
 B_0=1,\qquad
 B_n=-\frac1{n+1}\sum_{j=0}^{n-1}\binom{n+1}{j}B_j
 \quad(n\ge1)
\end{equation}
shows that every $B_n$ used here has denominator prime to $P$:
all indices are at most $1922<P-1$. Thus the entire calculation
can be performed in $\mathbb F_P$ without an undefined division.
Using $B_1=-1/2$ and $B_{2r+1}=0$ for $r\ge1$, the nine required
residues, in descending-index order, are
\begin{equation}\label{t10:Bvalues}
\begin{array}{c|rrrrrrrrr}
2r&1922&1920&1918&1916&1914&1912&1910&1908&1906\\\hline
 B_{2r}\bmod P&31528&45962&26974&45003&39580&2844&23974&18828&26318.
\end{array}
\end{equation}
Substitution into \eqref{t10:A} gives
\begin{equation}\label{t10:Amatrix}
\mathcal A\equiv
\begin{pmatrix}
25959&60986&10404&1273&2935\\
16142&62424&8365&2781&39264\\
13624&13203&159&63417&38163\\
51637&50007&47664&40198&52007\\
23028&22542&12445&19571&47569
\end{pmatrix}\pmod{65537}.
\end{equation}
Gaussian elimination without row exchanges has diagonal pivots
\[
25959,\quad2542,\quad6997,\quad61613,\quad26239.
\]
Their product is $60407$ modulo $65537$. The complete multiplier
in \eqref{t10:multiplier}, including its sign, is $49105$ modulo
$65537$. Therefore
\[
\det M_{5,962}\equiv49105\cdot60407\equiv15578\not\equiv0
     \pmod{65537},
\]
which proves \eqref{t10:main-certificate}.

\begin{proof}[Proof of Theorem~\ref{t10:complete}]
If $b\not\equiv962\pmod{2048}$, use the source order-ten theorem.
If $b\ge1920$, use the source effective eventual theorem with
$B_5=1920$. The only integer $b$ with
$6\le b<1920$ and $b\equiv962\pmod{2048}$ is $962$;
\eqref{t10:main-certificate} treats that value. These cases cover
all integers $b\ge6$. The inherited triangular transforms and
finite late-coordinate expansion then give the stated moment
reconstruction.\end{proof}

The supplied verifier \texttt{complete\_p10.py} implements
\eqref{t10:Brec}, \eqref{t10:Newton}, \eqref{t10:A}, and modular
Gaussian elimination, using only Python integer arithmetic. A
second verifier reconstructs the same matrix from the uncontracted
Pad\'e row sums and separately reduces exact rational Bernoulli
numbers obtained from SymPy. These checks agree entry by entry.
The prime is checked by trial division through its integer square
root. No floating-point determinant is used anywhere in the proof.
This is a finite arithmetic proof certificate, not a claim of
proof-assistant formalization.

\subsection{An ordinary convergent polygamma identity for every larger order}

The nonvanishing theorem makes the inherited reconstruction an
unconditional all-$q$ identity family at smaller order ten. Define,
for $0\le k<b$,
\begin{equation}\label{t10:weight}
 P_{b,k}(x)=\sum_{t=0}^{b-1}(M_{5,b}^{-1})_{kt}
 \left\{(x+\tfrac12)^{2b+2t+1}-(x-\tfrac12)^{2b+2t+1}\right\}
 =\sum_{r=0}^{2b-1}w_{b,k,r}x^{2r}.
\end{equation}
All coefficients are rational. Put
\[
 s_{p,0}=\frac1{p-1},\qquad
 s_{p,j}=\frac{B_{2j}(1/2)}{(2j)!}(p)_{2j-1}\quad(j\ge1),
 \qquad
 \alpha_j=\sum_{h=0}^js_{10,h}s_{2b,j-h},
\]
where $(p)_m$ is the rising factorial. Define the finite even polynomial
\begin{equation}\label{t10:subtraction}
 S_{b,k}(x)=\sum_{r=b+4}^{2b-1}w_{b,k,r}
          \sum_{j=0}^{r-b-4}\alpha_jx^{2(r-b-4-j)}.
\end{equation}
Then
\begin{equation}\label{t10:ordinary}
 \boxed{\sum_{n=0}^{\infty}
 \left\{P_{b,k}(x_n)\,
 \frac{\psi^{(9)}(n+1)\psi^{(2b-1)}(n+1)}{9!\,(2b-1)!}
                -S_{b,k}(x_n)\right\}=\zeta(2k).}
\end{equation}
Every sum in \eqref{t10:ordinary} converges absolutely. Indeed,
the centered Hurwitz expansion is
\[
 A_p(n)\sim\sum_{j\ge0}s_{p,j}x_n^{1-p-2j}.
\]
Thus \eqref{t10:subtraction} removes exactly every nondecaying
term of $P_{b,k}(x_n)A_{10}(n)A_{2b}(n)$ and leaves $O(x_n^{-2})$.
The polynomial difference in \eqref{t10:weight} is exactly the
inverse transform \eqref{t10:inverse}. Matrix inversion therefore
makes its continued sum $\zeta(2k)$. All removed monomials have
zero centered continued sum, since
$\zeta(-2j,1/2)=0$ for integers $j\ge0$, proving the ordinary
convergent identity. The polygamma substitution uses
$\psi^{(p-1)}(n+1)=(p-1)!\zeta(p,n+1)$ for even $p$.

The new contribution is the complete smaller-order-ten theorem,
including its single missing exact certificate. The centered
continuation, Pad\'e contraction, and conversion of reconstructed
coordinates into convergent sums retain their source attribution.
